\section{Graph Definitions}
\label{section:21-graph-definitions}

This section will introduce key notation and definitions used throughout the thesis. 

The following definitions describe the commonly agreed building blocks and descriptors of a graph in an abstract sense. 
\nrscomment{this wording is dangerously close to saying the primitives are agreed upon. maybe reword to say the abstract components and characteristics of graphs?}

\define{Graph ($G$)}{The data structures comprising of the related vertices and edges such that $G=\{V,E\}$. Also called a `network'. $|G|$ represents the size of the graph.}
\nrscomment{a graph is a singular data structure. size in number of nodes? edges? is it clearly defined?}

\define{Vertex ($V$)}{Represent the `entities' or `things'. Also called a node. $|V|$ represents the number of vertices or size of the vertex set.}
\define{Edge ($E$)}{Represent the relationships between vertices. Also called arcs or links. $|E|$  represents the number of vertices or the size of the edge set.}
\define{Attribute}{The labels, properties or other information that is encoded in the vertices or edges of a graph. Optional in graphs.}
\nrscomment{make the defn singular}

\define{Weight}{The numerical value assigned to an edge to describe its importance in a graph. Optional in graphs.}
\nrscomment{is it always importance? cant it represent some actual value in the world, like number of citations? instead of saying its optional, maybe say that graphs having edges with weights are termed weighted graphs.}

\define{Direction}{The flow of information across an edge. If a graph is directed, the order of source and destination vertex is important. If undirected, vertex order is not important.}
\nrscomment{just qualify the first sentence with '... across an edge, from source node to destination node. Graphs having edges with direction are termed directed graphs.' and then cut the second two sentences}


%The following definitions describe the aggregations, generalizations and descriptors of graphs: 
%
\nrscomment{id say graphs have a number of important characteristics, including... and then do the definitions}

\define{Topology}{The structure of a graph \nrscomment{missing 'that'?} is created when edges like \nrscomment{replace 'like' with 'link'?} vertices together \nrscomment{"...together, regardless of edge attributes or weights."}. It does not include attributes or weights. Extracting the latent descriptive power of the graph topology is the focus of most `classical' graph processing.}
\nrscomment{as it stands this isnt a definition, its a collection of words without a verb. im not sure the second sentence is really relevant to the definition. if you think it is, clarify what you mean by the latent descriptive power of the topology}

\define{Degree}{The count of edges at a vertex. Frequently split between degree$_{out}$ and degree$_{in}$ in directed graphs. Graphs are often described by their `average degree'. Larger average degrees indicate denser graphs.}
\define{Neighborhood}{The collection of vertices with edges directly connecting them to a given vertex. The $K$-hop neighborhood is the collection of vertices within $K$ edges of a given vertex (e.g. friend of a friend)}
\define{Diameter}{The shortest path between the two most distant vertices in a graph. Sparser graphs tend to have larger diameters.}
\nrscomment{why? need some explanation like there are more nodes between ... whatever. what is meant by distance? number of nodes traversed on the shortest path between?}


%The following definitions describe the characteristics of graphs: 

\define{Static}{A graph that is known apriori whose vertices, edges, weights and attributes do not change. Static graphs are generally simpler to process.}
\define{Dynamic}{A graph that evolves with vertices, edges, attributes and weights being created, updated and deleted over time. Also called a streaming graph. Dynamic graphs introduce temporal constraints on graph processing.}
\define{Heterogeneity}{Describes a graph with multiple `types' of vertex, for example, a social graph that has vertices to represent `person' and vertices to represent `organization'. Heterogeneous graphs are more complicated to process because of the requirement for type-checking \nrscomment{need a subject for type-checking to act on. nodes maybe.}.}
\define{Degree Distribution}{The distribution of vertex degrees in a given graph. Can follow any statistical distribution and correlates heavily with known graph processing bottlenecks.}
\define{Multi-graph}{A graph that allows multiple edges between the same two vertices. Common in knowledge graphs, and other graphs with multiple vertex and edge types.}
\define{Hyper-edge}{An edge that connects an arbitrary number of vertices}


\osullikomment{Graph Diagrams Here?}